\section{First fall comparisons and degree-four measurements}
\subsection{A comparison of two degree conventions}

Nagao's Definitions 5 and 6 distinguish a degree in
$\F_p[X_1,\ldots,X_N]$ from a degree after reduction modulo the field
equations. Lemma 4 compares $d_F(F)$ to $d'_F(F\cup S_{\rm fe})$, with
$S_{\rm fe}=\{X_i^p-X_i\}$ \cite{nagao2015}. The frozen extraction explicitly
places the field equations outside the true-side generator set in that
lemma. The programme's current certificate code makes the intended
comparison convention explicit \cite{nagaoCertificate}.

For the equality convention, a fall at $D$ has products $g_i f_i$ whose
maximum degree is exactly $D$, and whose nonzero sum has degree strictly
less than $D$. In the fake quantity every product and the sum is first
reduced modulo $S_{\rm fe}$. The requirement that the maximum is attained
must be checked; a low-degree element in a span alone is insufficient.

\begin{proposition}[Existing counterexample, equality convention]
Over $\F_2$, the universal inequality $d_F\leq d'_F$ fails. Two recorded
witnesses are
\[
 F_A=\{X^2Y,\;XY+X\},\qquad
 F_B=\{X_1X_2+X_3,\;X_1X_3+X_2\}.
\]
The current certificate tests assign fake degree $2$ and true degree $3$
to each witness, for either true-side placement of the field equations.
\end{proposition}

The witnesses and exact test assertions are in
\art{crypto-autoresearcher:tools/lemma4_inside_form.py} and its test module.
For $F_B$, take $g_1=1$ and $g_2=X_1+1$. Boolean reduction gives
\[
 \overline{f_1+(X_1+1)f_2}=X_2+X_3,
 \quad \deg\overline{f_1}=\deg\overline{(X_1+1)f_2}=2.
\]
The ordinary products attain degree $3$, not $2$. The tests separately
check all degree-two scalar combinations when field equations are
adjoined, and exhibit a degree-three fall. These checks are finite
certificates for the stated polynomial systems, not elliptic-curve
instances or measurements of an F4 run.

\paragraph{The printed inequality reading is a separate axis.}
The source prints a $\geq$ form of the degree condition. Under the
programme's padding analysis that reading collapses the quantity to one
plus the least degree of a nonzero ideal element. The current test class
\art{TestTheGeReadingVoidsOnlyOnePlacement} distinguishes the true-side
placement as follows, for both witnesses:
\begin{center}
\begin{tabular}{llll}
\toprule
Condition reading & True-side generators & True & Fake \\
\midrule
Equality & $F$ & 3 & 2 \\
Equality & $F\cup S_{\rm fe}$ & 3 & 2 \\
Literal $\geq$ & $F$ & 3 & 2 \\
Literal $\geq$ & $F\cup S_{\rm fe}$ & 2 & 2 \\
\bottomrule
\end{tabular}
\end{center}
Thus the assertion that the literal reading always voids the
counterexample is too broad. The historical finding
\art{KN-FIND-936151} contains that broader wording; the later tests
explicitly correct it. Both are retained in the source inventory. The
frozen extracted Lemma 4 names $F$ on the true side, but a final
publication still owes a fresh PDF-level transcription of both definitions
and their intended reading. These sources do not prove or disprove
Nagao's Proposition 5 by another argument, nor a complexity conclusion
for ECDLP. The programme itself withholds the broader derived claim.

\subsection{What the binary degree measurements actually establish}

Semaev's Assumption 1 concerns $d_{F4}$ for the Boolean system equivalent
to his chained equations, with subspace dimension $k=\lceil n/m\rceil$
\cite{semaev2015}. It is distinct from Nagao's separately numbered
Assumption 1, from first fall degree, from a one-pass Macaulay rank, and
from a bounded-degree mutant closure. Confusing these quantities can
turn a valid computation into an invalid refutation.

The closure-at-scale study \cite{closureScale} reports eight unsatisfiable
$m=3$ draws at $(n,\ell)=(19,7)$ and $(21,7)$ for which the complete
degree-four mutant closure saturates without containing $1$. It also
reports degree-four refutation on its other listed cells, including the
$m=4$ cells through $\ell=4$. This is a failure of a uniform
constant-four closure model in that representation. It does not follow
merely from a timeout: the report says these particular closures
converged. A new verifier would still need the retained instance, closure
basis, closure-stability checks and non-membership evidence before
certifying an individual negative instance afresh.

\status{The broad claim that these records refute Semaev's Assumption 1
is not established.} The programme's own correction
\art{KN-FIND-9cba50} records that an earlier ladder used
Frobenius-invariant subspaces of dimension $\operatorname{ord}_n(2)$,
which need not equal $\lceil n/m\rceil$, and that its degree meters were
not interchangeable with $d_{F4}$. Conversely, the two $m=3$ cells just
listed do have $\ell=\lceil n/m\rceil$; calling them off-diagonal does
not fix the metric issue. A verified bridge from saturated mutant
closure to the precise F4 quantity remains necessary in either case.

There is also positive counterevidence against a naive degree-meter
identification. At the separately archived $n=17,m=2,\ell=9$ cell,
plain degree-four Macaulay left 62 unsatisfiable targets unresolved;
the subsequent mutant closure refuted all 62 with checked witnesses
\cite{closureCorrection}. Therefore non-refutation by plain Macaulay
was not itself evidence that mutant refutation degree exceeded four.
No statement about cryptographic-size degrees follows from this finite
example. Likewise, fitting either of two surviving degree-growth models
does not logically refute every possible $m=4$ degree surface outside
the measured cells.
